% =============================================================================
% 6.4 ESTRATEGIA DE INDEXACIÓN Y OPTIMIZACIÓN
% =============================================================================
\section{Estrategia de Indexación y Optimización de Consultas}
\label{sec:indexacion}

El esquema define alrededor de \textbf{80 índices} (incluyendo los implícitos de PK y
UNIQUE). La estrategia no es indexar indiscriminadamente, sino atacar los patrones de
consulta reales del backend: filtrado por dueño (\texttt{profile\_id}), exclusión de
borrados lógicos y ordenación por relevancia o fecha.

% -----------------------------------------------------------------------------
\subsection{Índices B-Tree sobre identificadores y llaves foráneas}
\label{subsec:btree}

Todas las llaves foráneas de uso frecuente cuentan con un índice B-Tree explícito que
acelera los \textit{joins} y los filtros por dueño. La tabla~\ref{tab:tipos-indice}
clasifica los índices del esquema por tipo y propósito.

\vspace{0.2cm}
\noindent
\renewcommand{\arraystretch}{1.3}
\fontsize{8.5}{10.2}\selectfont\sffamily
\begin{tabularx}{\textwidth}{|>{\bfseries\color{colorPrimario}}l|c|X|}
\hline
\rowcolor{headerTabla}
\textbf{Tipo de índice} & \textbf{Aprox.} & \textbf{Propósito y ejemplo} \\
\hline
B-Tree simple (FK) & $\sim$20 & Acelera \textit{joins} y filtros por dueño. Ej.: \texttt{idx\_projects\_profile\_id}. \\
\hline
\rowcolor{grisTecnico}
Compuesto & $\sim$12 & Cubre filtro + orden en una pasada. Ej.: \texttt{idx\_notifications\_profile (profile\_id, is\_read, created\_at DESC)}. \\
\hline
Parcial (\textit{soft-delete}) & $\sim$10 & Indexa solo filas activas (\texttt{WHERE deleted\_at IS NULL}), reduciendo tamaño. \\
\hline
\rowcolor{grisTecnico}
Único parcial & $\sim$6 & Unicidad condicional. Ej.: \texttt{uq\_profile\_tag} único por \texttt{(profile\_id, tag\_id)} solo si activo. \\
\hline
Funcional & 1 & \texttt{idx\_login\_attempts\_email} sobre \texttt{lower(email)} para búsqueda case-insensitive. \\
\hline
\end{tabularx}
\label{tab:tipos-indice}

\subsection{Índices compuestos y parciales para búsquedas optimizadas}
\label{subsec:compuestos}

Los \textbf{índices parciales} son la pieza central de la estrategia de
\textit{soft-delete}: como la mayoría de las consultas filtran por
\texttt{deleted\_at IS NULL}, indexar solo las filas activas reduce el tamaño del índice
y mejora la selectividad. La tabla~\ref{tab:indices-clave} muestra ejemplos
representativos extraídos del esquema real.

\vspace{0.2cm}
\noindent
\renewcommand{\arraystretch}{1.25}
\fontsize{8}{9.6}\selectfont\sffamily
\begin{tabularx}{\textwidth}{|>{\ttfamily\color{colorPrimario}\scriptsize}l|X|}
\hline
\rowcolor{headerTabla}
\sffamily\textbf{Índice} & \sffamily\textbf{Definición y motivación} \\
\hline
idx\_projects\_profile\_active & \texttt{(profile\_id, is\_featured DESC, created\_at DESC) WHERE deleted\_at IS NULL}. Lista los proyectos activos de un perfil, destacados primero. \\
\hline
\rowcolor{grisTecnico}
idx\_work\_experiences\_active & \texttt{(profile\_id, is\_current DESC, start\_date DESC) WHERE deleted\_at IS NULL}. Línea de tiempo laboral. \\
\hline
idx\_notifications\_profile & \texttt{(profile\_id, is\_read, created\_at DESC)}. Bandeja de notificaciones del usuario. \\
\hline
\rowcolor{grisTecnico}
idx\_profile\_likes\_company\_stage & \texttt{(company\_profile\_id, stage, stage\_order)}. Render del pipeline Kanban por fase. \\
\hline
uq\_chat\_messages\_client\_id & \texttt{(chat\_id, client\_message\_id) WHERE client\_message\_id IS NOT NULL}. Garantiza idempotencia de mensajes. \\
\hline
\rowcolor{grisTecnico}
idx\_portfolio\_settings\_slug & \texttt{(slug) WHERE slug IS NOT NULL AND slug <> ''}. Resolución de portafolios públicos por slug. \\
\hline
idx\_psc\_active\_lookup & \texttt{(profile\_id, purpose, code) WHERE used = false}. Validación rápida de OTP vigentes. \\
\hline
\rowcolor{grisTecnico}
idx\_tss\_niche\_week & \texttt{(niche, week\_start, rank\_position)}. Cálculo del leaderboard por nicho. \\
\hline
\end{tabularx}
\label{tab:indices-clave}

\begin{NotaTecnica}
El índice \comando{uq\_chat\_messages\_client\_id} merece especial atención: al ser un
\textbf{índice único parcial} sobre \texttt{client\_message\_id}, implementa la
idempotencia de mensajes a nivel de motor. Si el cliente reintenta un envío con el mismo
\texttt{client\_message\_id}, el \texttt{INSERT} viola la unicidad y se evita el
duplicado, sin necesidad de lógica adicional en la aplicación.
\end{NotaTecnica}
